\chapter{Bundles and connections at a cut}\label{ch:02}

Identifying two spacetime points does not tell us which vectors in their two fibers represent the same field value. For a charged field, even the statement that two complex numbers agree presupposes a comparison of frames. We will construct the fiber identification first. It will tell us how to compare matter fields and connections, and it will leave a measurable choice on a spatial circle.

\section{Lifting a face map}

Let $G$ be a Lie group, possibly disconnected. A principal $G$-bundle $P\to R$ has a right $G$-action, free and transitive on each fiber, and smooth local product charts. A local section $e_i$ chooses a frame: each point in the fiber has a unique expression $e_i(x)g$. An associated field with typical fiber $W$ and left $G$-action takes values in $P\times_G W$, with $[p g,w]=[p,g w]$. This construction works for a vector representation, but does not require $W$ to be linear.

For the face map $f:F_-\to F_+$ of Chapter~\ref{ch:01}, seek an equivariant isomorphism $D:P_-|_{F_-}\to P_+|_{F_+}$ covering $f$. Thus the test is whether $P_-|_{F_-}$ and $f^*(P_+|_{F_+})$ are isomorphic, not whether each can be trivialized on a small patch. To carry it out, choose frames on corresponding patches and write
\[
 e_j^\pm=e_i^\pm g_{ij}^\pm,\qquad
 D(e_i^-(y))=e_i^+(f(y))h_i(y).
\]
Evaluating $D(e_j^-)$ first in the $i$ frame and then in the $j$ frame gives
\begin{equation}
 h_i g_{ij}^-=(g_{ij}^+\circ f)h_j .\label{eq:02-clutch}
\end{equation}
Conversely, smooth $G$-valued $h_i$ satisfying \eqref{eq:02-clutch} define $D$ by $e_i^- g\mapsto e_i^+ h_i g$ at corresponding points. The equation is exactly what makes the definition independent of $i$ \cite{Steenrod}. Under a passive frame change $e_i^\pm\mapsto e_i^\pm k_i^\pm$ the same $D$ has matrix $(k_i^+\circ f)^{-1}h_i k_i^-$. No bundle has changed in this calculation.

Several bundles can be present. If a map $b_\pm:E_\pm\to H_\pm$ is part of the structure, the two proposed lifts must satisfy $D_H b_-=b_+D_E$. For example, an identification of spin bundles must project to the specified identification of oriented orthonormal frames; an abstract isomorphism of the two total spaces is not enough. The relevant frame comparison acts on the full spacetime tangent space, including the transverse direction identified by the signed collars. In contrast, an associated bundle has an induced lift $[p,w]\mapsto[D(p),w]$ by equivariance. There is no second independent twist for each associated species unless the theory actually supplies such structure.

\section{Smooth clutching and the connection law}

Construct $P$ over the sewn $M$ by identifying $p$ on the seam with $D(p)$. To see smoothness, choose product descriptions of the two collar bundles as pullbacks of their face bundles. Parallel transport along the collar lines with any auxiliary connection gives such descriptions for an ordinary principal bundle. This is a choice of charts, not of the physical connection. In extended frames let $g_-,g_+$ be the group coordinates on each side. A common fiber coordinate over the signed base box is $h_i(y)g_-$ on $s\leq0$ and $g_+$ on $s\geq0$. The inverse uses $h_i(y)^{-1}$ on the minus sector. Equation~\eqref{eq:02-clutch} makes the transition to the next product frame a single smooth map. The right action is ordinary right multiplication throughout this box. Thus the quotient is locally $B\times G$ with principal fibers, not merely a set of identified field components.

At a joint choose compatible product descriptions for the whole diagram, or give actual smooth signed transitions realizing the chosen lifts. Write $D_{ij}$ for transport from frame $j$ to frame $i$. On a common triple overlap one needs $D_{ij}D_{jk}=D_{ik}$ and $D_{ji}=D_{ij}^{-1}$. When only adjacent sectors have been paired at a four-sector vertex, the two routes around the square from one sector to the opposite one must agree there, with compatible smooth transitions. Otherwise going around a single point could identify $p$ with $p g$ for $g\neq 1$, destroying freeness. These are local conditions at the \emph{same} point, not an instruction to set the transport around a noncontractible spatial loop to one. Rigid structures, such as a prescribed framing or a fixed connection, require their own compatibility in the collar; an auxiliary product description need not preserve them.

A principal connection specifies which directions in $P$ are horizontal, so a base path can be lifted by parallel transport. Its Lie-algebra-valued form $\omega$ vanishes on horizontal vectors, returns $X$ on the vertical vector generated by $X$, and obeys $R_g^*\omega=\operatorname{Ad}_{g^{-1}}\omega$ under right multiplication. Its local potential is $A=e^*\omega$. Differentiating a change of section $e\mapsto eh$ gives $(eh)^*\omega=\operatorname{Ad}_{h^{-1}}(e^*\omega)+h^{-1}dh$: the first term differentiates $e$, the second is the vertical motion of $h$. In our common signed chart $e_-=e_+h$, with base pullbacks understood. Solving for $A_+$ gives
\[
 A_+=\operatorname{Ad}_h A_--(dh)h^{-1} .
\]
In a vector representation this is just $(d+A_+)(h z_-)=h(d+A_-)z_-$ for every matter section with components $z_+=h z_-$. The connection law is affine, not tensorial: only the difference of two connections transforms homogeneously. In particular a $U(1)$ change $h=e^{i c s}$ gives $A_+=-i c\,ds$ if $A_-=0$, even though $h=1$ on the face. The displayed law relates two descriptions of \emph{one} connection. To test two independently given regional connections, transform their one-sided coefficients into one signed frame, including the $ds$ component and the derivatives of the realized frame change, and compare the seam jets there. Equality of tangential pullbacks alone is weaker. We make the smoothness criterion precise in Chapter~\ref{ch:03}.

\section{Which changes identify the same sewing?}

A marked sewing remembers the labelled regions, the chosen cut and the prescribed exterior data. An isomorphism may reparametrize a labelled region by an allowed automorphism of the \emph{whole} region; it does not erase the decomposition. Specify which exterior fields, sources and frames are fixed, and which may be transported. Let $A_a$ consist of the actual automorphisms of region $a$ respecting these choices and the fixed structure. A symmetry defined only on a cut need not extend into the region subject to those requirements.

For a realized smooth sewing, write $D$ also for its full seam transition data (base and bundle maps with their signed Taylor coefficients). Write $\widehat u_a$ for the corresponding restriction of $u_a\in A_a$. An isomorphism from data $D$ to data $D'$ satisfies
\begin{equation}
 \widehat u_+ D=D'\widehat u_- .\label{eq:02-mark}
\end{equation}
Indeed a point transported across the seam and then reparametrized must agree with one reparametrized first and transported across the new seam. The same equality on derivatives is what makes the piecewise map smooth; the descent argument of the next chapter proves this for maps as well as fields. Conversely a smooth isomorphism preserving the marks restricts to actual regional maps satisfying \eqref{eq:02-mark}. Thus the classes are the orbits of admissible $D$ under $D\mapsto\widehat u_+D\widehat u_-^{-1}$. The stabilizer of $D$ consists of the \emph{pairs} $(u_-,u_+)$ satisfying \eqref{eq:02-mark} with $D'=D$. For instance, allowed maps supported away from the cut can be invisible in $\widehat u_a$ and still be distinct automorphisms. Several faces of one region use restrictions of the \emph{same} map. This describes marked sewings; an isomorphism allowed to move or forget the cut need not restrict to such maps.

\section{Worked check: two signs on a cylinder}

Build $\mathbb R_t\times S^1$ from two connected spatial strips, each carrying the trivial principal $G=\{+1,-1\}$ bundle. Fix the base and impose no exterior framing. In order around the circle let the seam lifts be $a$ from strip 1 to strip 2 and $b$ from strip 2 back to strip 1. A smooth vertical automorphism on a connected strip is a \emph{single constant} $u_i=\pm1$ because the group is discrete. Thus
\[
 a'=u_2 a u_1^{-1},\qquad b'=u_1 b u_2^{-1},\qquad b'a'=ba .
\]
Lift a spatial circuit at fixed $t$, starting with fiber element $1$. It stays constant in each strip frame, jumps to $a$ at the first seam and returns as $ba$ at the second. Its holonomy $H=ba$ therefore has a direct path-lifting meaning. Choosing $u_1=1,u_2=a^{-1}$ puts the pair in the form $(1,H)$. There are exactly two classes, each with diagonal stabilizer $u_1=u_2$. In the associated sign representation a field on the covering line obeys $\phi(x+L)=H\phi(x)$ for circumference $L$. Allowing an unrelated sign at each end of a strip would incorrectly erase the antiperiodic possibility. This is the holonomy of a discrete covering bundle; in a connected group a connection can have nontrivial holonomy even on a topologically trivial bundle.

\begin{exercise}
On the boundary of a disk prescribe a $U(1)$ gauge transformation $v(\theta)=e^{im\theta}$, $m\in\mathbb Z$. For which $m$ does it extend to a smooth gauge transformation on the disk? What does this imply for the winding of a transition between two trivial disk bundles?

\emph{Hint.} For an extension $u$, the form $u^{-1}du$ on the disk is closed. Stokes' theorem gives $0=\oint v^{-1}dv=2\pi i m$. The zero case extends by the constant map. Frame changes from either disk have zero boundary winding and cannot change the winding of the transition.
\end{exercise}
